\section{Scalar deep-linear hierarchy and failure beyond two factors}
\label{sec:deep-linear}

Consider the scalar width-one depth-$d$ linear network
\begin{equation}
  f(w)=\prod_{i=1}^d w_i
  \label{eq:scalar-network}
\end{equation}
with scalar residual $r=\ell'(f)$.  On a fixed sign component, put
\begin{equation}
  q_i=w_i^2>0,
  \qquad
  e_k(q)=\sum_{1\leq i_1<\cdots<i_k\leq d}q_{i_1}\cdots q_{i_k},
  \qquad e_0=1.
\end{equation}
Then $f^2=e_d$ and
\begin{equation}
  K=\sum_{i=1}^d\left(\frac{\partial f}{\partial w_i}\right)^2
  =e_{d-1}(q).
  \label{eq:scalar-K}
\end{equation}
Gradient flow gives the common translation law
\begin{equation}
  \dot q_i=-2fr
  \qquad \text{for every }i.
  \label{eq:q-common-velocity}
\end{equation}

\begin{proposition}[Finite symmetric-polynomial hierarchy]
\label[proposition]{prop:symmetric-hierarchy}
For $k=1,\ldots,d$,
\begin{equation}
  \boxed{\dot e_k=-2fr(d-k+1)e_{k-1}.}
  \label{eq:hierarchy}
\end{equation}
In particular,
\begin{equation}
  \dot f=-e_{d-1}r=-Kr,
  \qquad
  \boxed{\dot K=-4fr e_{d-2}.}
  \label{eq:K-hierarchy}
\end{equation}
Thus $(f,e_1,\ldots,e_{d-1})$ is an exact finite-dimensional state.
\end{proposition}

\begin{proof}
Since every $q_i$ has velocity $-2fr$,
\begin{equation}
  \dot e_k=-2fr\sum_{i=1}^d\frac{\partial e_k}{\partial q_i}.
\end{equation}
Every monomial of degree $k-1$ occurs in exactly $d-k+1$ of the derivatives, so
\begin{equation}
  \sum_i\frac{\partial e_k}{\partial q_i}=(d-k+1)e_{k-1}.
\end{equation}
The stated equations follow.
\end{proof}

At $d=1$, $K\equiv1$.  At $d=2$, \eqref{eq:K-hierarchy} gives
\begin{equation}
  \dot f=-Kr,
  \qquad
  \dot K=-4fr,
  \label{eq:scalar-depth-two}
\end{equation}
which is the scalar specialization of \cref{prop:two-factor-closure}.  From $d=3$ onward, $e_{d-2}$ is a new coordinate.  The next theorem shows directly that it is not determined by $(f,K)$ on generic fibers.

\begin{theorem}[Regular-fiber obstruction from three factors onward]
\label[theorem]{thm:fiber-obstruction}
Let $d\geq3$ and define
\begin{equation}
  \pi_d(q)=\bigl(e_d(q),e_{d-1}(q)\bigr),
  \qquad q\in\R_{>0}^d.
\end{equation}
At every point for which the reciprocals $u_i=1/q_i$ take at least three distinct values,
\begin{equation}
  \rank D\bigl(e_d,e_{d-1},e_{d-2}\bigr)=3.
  \label{eq:rank-three}
\end{equation}
Consequently, the fiber of $\pi_d$ is locally a smooth manifold of dimension $d-2$ containing curves along which $e_{d-2}$ is nonconstant.  For $d=3$ the fiber itself is locally a curve, and $e_1$ varies along it.

Therefore, for any loss and output value with $fr\neq0$, there is no leaf-independent single-valued law
\begin{equation}
  \dot K=\mathcal F(f,K,r)
  \label{eq:no-closure}
\end{equation}
on the full positive state space.  Three factors are the first scalar linear architecture at which the universal two-factor closure fails.
\end{theorem}

\begin{proof}
The change of variables $u_i=1/q_i$ is a diffeomorphism of the positive orthant.  Put
\begin{equation}
  P=e_d(q),
  \qquad
  S=\frac{e_{d-1}(q)}{e_d(q)}=\sum_i u_i,
  \qquad
  T=\frac{e_{d-2}(q)}{e_d(q)}=\sum_{i<j}u_i u_j.
  \label{eq:reciprocal-coordinates}
\end{equation}
The transformations $(e_d,e_{d-1},e_{d-2})\leftrightarrow(\log P,S,T)$ are locally invertible for $P>0$, so it suffices to study the differentials in $u$-coordinates:
\begin{equation}
  \dd\log P=-\sum_i\frac{\dd u_i}{u_i},
  \qquad
  \dd S=\sum_i\dd u_i,
  \qquad
  \dd T=\sum_i(S-u_i)\dd u_i.
  \label{eq:three-differentials}
\end{equation}
The first two forms are independent unless all $u_i$ are equal.  Suppose that $\dd T=A\,\dd S+B\,\dd\log P$.  Comparing the coefficient of $\dd u_i$ gives
\begin{equation}
  S-u_i=A-\frac{B}{u_i}.
\end{equation}
Hence every $u_i$ is a root of the same quadratic
\begin{equation}
  z^2-(S-A)z-B=0.
\end{equation}
This is impossible when at least three distinct reciprocal values occur.  The three differentials are therefore independent, proving \eqref{eq:rank-three}.

The constant-rank theorem now makes the level set of $(e_d,e_{d-1})$ a local manifold of dimension $d-2$.  Since $\dd e_{d-2}$ is not in the span of $\dd e_d$ and $\dd e_{d-1}$, it has a nonzero derivative in some tangent direction to that fiber; a smooth curve tangent to that direction keeps $(e_d,e_{d-1})$ fixed while changing $e_{d-2}$.  Along the positive sign component, fixed $e_d$ means fixed $f$, and fixed $e_{d-1}$ means fixed $K$.  Equation \eqref{eq:K-hierarchy} then gives different values of $\dot K$ whenever $fr\neq0$.
\end{proof}

\begin{remark*}[Explicit depth-three witness]
For $d=3$, take
\[
  q^A=\left(\frac25,\frac54,2\right),
  \qquad
  q^B=\left(\frac59,\frac35,3\right).
\]
Then $e_3(q^A)=e_3(q^B)=1$ and $e_2(q^A)=e_2(q^B)=19/5$, whereas
\[
  e_1(q^A)=\frac{73}{20},
  \qquad
  e_1(q^B)=\frac{187}{45}.
\]
Thus the two states have the same positive output $f=1$ and the same tangent kernel $K=19/5$, but \eqref{eq:K-hierarchy} gives different kernel velocities whenever $r(1)\neq0$.
\end{remark*}

The condition of three distinct reciprocal coordinates is open and dense.  The obstruction is therefore not an isolated pair of parameter states but a generic local property of the positive deep-linear state space.  It rules out a \emph{universal} closure on that space.  It does not rule out closure on a selected invariant leaf or on a more restrictive set of trajectories.

\subsection{Exact closure on each balancing leaf}

Equation \eqref{eq:q-common-velocity} implies that
\begin{equation}
  c_i=q_i-q_d,
  \qquad i=1,\ldots,d-1,
  \qquad \dot c_i=0.
  \label{eq:leaf-labels}
\end{equation}
Fix a positive leaf $c$ and write $q_d=s$, $q_i=s+c_i$.  Define
\begin{equation}
  \kappa_c(s)=e_{d-1}(s+c_1,\ldots,s+c_{d-1},s).
\end{equation}

\begin{proposition}[Leaf-conditioned exact closure]
\label[proposition]{prop:leaf-closure}
On every interval where all $q_i$ are positive,
\begin{equation}
  \kappa_c'(s)=2e_{d-2}(q(s,c))>0.
  \label{eq:kappa-derivative}
\end{equation}
Hence $s=\kappa_c^{-1}(K)$ and
\begin{equation}
  \boxed{
  \dot K=-4fr\,e_{d-2}\bigl(q(\kappa_c^{-1}(K),c)\bigr),
  \qquad \dot c=0.}
  \label{eq:leaf-closure}
\end{equation}
Thus the failure in \cref{thm:fiber-obstruction} is transverse to the balancing foliation: a conserved label, not temporal history, restores exact closure.
\end{proposition}

\begin{proof}
Differentiation under a common translation and \cref{prop:symmetric-hierarchy} give
\begin{equation}
  \frac{\dd}{\dd s}e_{d-1}(q(s,c))=2e_{d-2}(q(s,c)).
\end{equation}
Positivity makes this derivative strictly positive.  Inversion and \eqref{eq:K-hierarchy} yield \eqref{eq:leaf-closure}.
\end{proof}

\subsection{Depth three in invariant coordinates}

For $d=3$, let
\begin{equation}
  I(q)=\frac12\sum_{i=1}^3(q_i-\bar q)^2,
  \qquad
  \bar q=\frac13e_1(q).
  \label{eq:I}
\end{equation}
The common translation \eqref{eq:q-common-velocity} gives $\dot I=0$, and
\begin{equation}
  I=\frac{e_1^2}{3}-e_2.
  \label{eq:I-identity}
\end{equation}
Since $K=e_2$ and $q_i>0$,
\begin{equation}
  e_1=\sqrt{3(K+I)}.
\end{equation}
The exact closed system is therefore
\begin{equation}
  \boxed{
  \dot f=-Kr,
  \qquad
  \dot K=-4fr\sqrt{3(K+I)},
  \qquad
  \dot I=0.}
  \label{eq:depth-three-closure}
\end{equation}
Because $f^2=e_3$, $K=e_2$, and $I$ recovers $e_1$, the state $(f,K,I)$ is simply the complete symmetric state $(e_3,e_2,e_1)$ written in invariant coordinates.  The one-scalar augmentation is not a further compression of the hierarchy.
